\chapter{Quantum Angular Momentum}

\sourcepage{1}
\section{Angular-momentum operators}
Recall the orbital angular-momentum operators
\begin{align*}
  \widehat L_x
  &=y\widehat p_z-z\widehat p_y
  =-\ii\hbar\left(y\frac{\partial}{\partial z}-z\frac{\partial}{\partial y}\right),\\
  \widehat L_y
  &=z\widehat p_x-x\widehat p_z
  =-\ii\hbar\left(z\frac{\partial}{\partial x}-x\frac{\partial}{\partial z}\right),\\
  \widehat L_z
  &=x\widehat p_y-y\widehat p_x
  =-\ii\hbar\left(x\frac{\partial}{\partial y}-y\frac{\partial}{\partial x}\right).
\end{align*}
Equivalently,
\[
  \widehat{\bm L}=-\ii\hbar\,\bm r\times\nabla.
\]
The commutation relations are
\[
  [\widehat L_y,\widehat L_z]=\ii\hbar\widehat L_x,
  \qquad
  [\widehat L_z,\widehat L_x]=\ii\hbar\widehat L_y,
  \qquad
  [\widehat L_x,\widehat L_y]=\ii\hbar\widehat L_z,
\]
or symbolically
\[
  \ii\hbar\widehat{\bm L}=\widehat{\bm L}\times\widehat{\bm L}.
\]
Also,
\[
  [\widehat L^2,\widehat L_z]=0
\]
(and similarly for $x$ and $y$ components).

For example, if $L^2=30\hbar^2$, then $\ell=5$, and the allowed $L_z$ eigenvalues are
\[
  L_z\in\{-5\hbar,-4\hbar,\ldots,5\hbar\}.
\]
The notation used later is
\[
  \widehat{\bm J}:\text{ total angular momentum},
  \qquad
  \widehat{\bm L}:\text{ orbital angular momentum},
  \qquad
  \widehat{\bm S}:\text{ spin}.
\]

\sourcepage{2}
\section{Eigenvalues and spherical harmonics}
The simultaneous eigenvalues of $\widehat J^2$ and $\widehat J_z$ are
\[
  J^2=\hbar^2j(j+1),
  \qquad
  J_z=\hbar m_j,
  \qquad
  m_j=-j,-j+1,\ldots,j.
\]
The eigenfunctions of $\widehat L^2$ and $\widehat L_z$ are the spherical harmonics,
\begin{equation}
  Y_\ell^m(\theta,\phi)
  =\left[\frac{2\ell+1}{4\pi}
  \frac{(\ell-m)!}{(\ell+m)!}\right]^{1/2}
  P_\ell^m(\cos\theta)\ee^{\ii m\phi}.
  \label{eq:qam-spherical-harmonics}
\end{equation}
An arbitrary angular function can be expanded as
\[
  f(\theta,\phi)
  =\sum_{\ell=0}^{\infty}\sum_{m=-\ell}^{\ell}
  a_{\ell m}Y_\ell^m(\theta,\phi),
\]
with
\[
  a_{\ell m}
  =\int_0^{2\pi}\dd\phi\int_{-1}^{1}\dd\mu\,
  Y_\ell^{m*}(\mu,\phi)f(\mu,\phi).
\]

\section{Addition of angular momenta}
For
\[
  \widehat{\bm L}=\widehat{\bm L}_1+\widehat{\bm L}_2,
\]
we have
\[
  \widehat L^2=(\widehat{\bm L}_1+\widehat{\bm L}_2)^2
  =\widehat L_1^2+\widehat L_2^2+2\widehat{\bm L}_1\cdot\widehat{\bm L}_2,
\]
and
\[
  \widehat L_z=\widehat L_{1z}+\widehat L_{2z}.
\]
The uncoupled product state is denoted
\[
  \ketc{\ell_1,\ell_2,m_1,m_2}.
\]

\sourcepage{3}
For fixed $\ell_1,\ell_2,m_1,m_2$, the uncoupled state
\[
  \ketc{\ell_1,\ell_2,m_1,m_2}
\]
is characterized by the mutually commuting operators $\widehat L_1^2$, $\widehat L_{1z}$, $\widehat L_2^2$, and $\widehat L_{2z}$.

An equally valid description is the coupled state
\[
  \ketc{\ell,m,\ell_1,\ell_2},
\]
characterized by the mutually commuting operators $\widehat L^2$, $\widehat L_z$, $\widehat L_1^2$, and $\widehat L_2^2$.

Thus there are two sets of ``good quantum numbers'':
\begin{center}
\begin{tabular}{c c}
\toprule
Coupled & Uncoupled\\
\midrule
$\ell,m,\ell_1,\ell_2$ & $\ell_1,\ell_2,m_1,m_2$\\
\bottomrule
\end{tabular}
\end{center}
The sketches emphasize the vector-coupling picture: in the coupled representation, $\bm L_1$ and $\bm L_2$ add to a fixed resultant $\bm L$; in the uncoupled representation, their separate projections are specified.

\sourcepage{4}
\section{Coupled representation}
The state $\ketc{\ell,m,\ell_1,\ell_2}$ is a simultaneous eigenstate of
\[
  \widehat L^2
  =\widehat L_1^2+\widehat L_2^2+2\widehat{\bm L}_1\cdot\widehat{\bm L}_2,
  \qquad
  \widehat L_z=\widehat L_{1z}+\widehat L_{2z},
  \qquad
  \widehat L_1^2,
  \qquad
  \widehat L_2^2.
\]
It can be expressed in the uncoupled product basis,
\[
  \ketc{\ell,m,\ell_1,\ell_2}
  =f(\Omega_1,\Omega_2)
  =\sum_{m_1=-\ell_1}^{\ell_1}
   \sum_{m_2=-\ell_2}^{\ell_2}
  f_{\ell m m_1m_2}
  Y_{\ell_1}^{m_1}(\Omega_1)Y_{\ell_2}^{m_2}(\Omega_2).
\]
The expansion coefficients are obtained by projection onto the product harmonics.  Conservation of angular momentum restricts the allowed total quantum numbers to
\[
  |\ell_1-\ell_2|\le\ell\le\ell_1+\ell_2,
  \qquad
  m=m_1+m_2.
\]

\sourcepage{5}
Accordingly, the sum need include only pairs satisfying $m_1+m_2=m$:
\begin{equation}
  \ketc{\ell,m,\ell_1,\ell_2}
  =\sum_{m_1+m_2=m}
  f_{\ell m m_1m_2}
  Y_{\ell_1}^{m_1}Y_{\ell_2}^{m_2}.
  \label{eq:qam-coupled-expansion}
\end{equation}
Equivalently,
\begin{equation}
  \ketc{\ell,m,\ell_1,\ell_2}
  =\sum_{m_1+m_2=m}
  C_{m_1m_2m}^{\ell_1\ell_2\ell}
  \ketc{\ell_1,\ell_2,m_1,m_2},
  \label{eq:qam-cg-expansion}
\end{equation}
where
\[
  C_{m_1m_2m}^{\ell_1\ell_2\ell}
  =\brac{\ell_1m_1;\ell_2m_2\mid\ell m}
\]
is a Clebsch--Gordan coefficient.

The interpretation is that
\[
  \left|C_{m_1m_2m}^{\ell_1\ell_2\ell}\right|^2
\]
is the probability that, in a state with fixed $\ell,m,\ell_1,\ell_2$, measurements of the component projections give $m_1$ and $m_2$.

For example, expand $\ketc{1,-1,1,1}$ in the uncoupled basis.  Since $m=-1$, the allowed pairs are $(m_1,m_2)=(0,-1)$ and $(-1,0)$:
\[
  \ketc{1,-1,1,1}
  =\frac{1}{\sqrt2}\ketc{1,1,0,-1}
  -\frac{1}{\sqrt2}\ketc{1,1,-1,0},
\]
with the signs following the convention used in the source notes.
